% Faithful textbook-style TikZ reconstruction of the official 2023 IMO Problem 5 example.
\begin{center}
\begin{tikzpicture}[x=1cm,y=1cm,line cap=round,line join=round]
    \def\circleRadius{0.235}
    \foreach \row in {1,...,6}{
        \foreach \column in {1,...,\row}{
            \coordinate (r\row c\column) at ({(\column-(\row+1)/2)*0.50},{(6-\row)*0.44});
        }
    }
    \foreach \row in {1,...,6}{
        \foreach \column in {1,...,\row}{
            \draw[line width=0.45pt,fill=white] (r\row c\column) circle (\circleRadius);
        }
    }
    % Exactly one red circle in each row, as in the official example.
    \foreach \r/\c in {1/1,2/1,3/3,4/3,5/4,6/1}{
        \draw[line width=0.45pt,fill=exampleRed] (r\r c\c) circle (\circleRadius);
    }
    % A ninja path through two red circles (rows 1 and 4).
    \draw[line width=1.15pt] (r1c1.center)--(r2c2.center)--(r3c2.center)
        --(r4c3.center)--(r5c3.center)--(r6c3.center);
    \draw[latex-latex,line width=0.45pt] (-1.50,-0.37)--(1.50,-0.37);
    \node[below] at (0,-0.39) {$n=6$};
\end{tikzpicture}
\end{center}
